\documentclass[../../main.tex]{subfiles}% 注意这里的文件路径不能用 ./main.tex ,否则用latexmk编译子文件会报错
\graphicspath{{\subfix{./image/}}{\subfix{../../image/}}} % 指定图片引用路径,后续可以直接使用图片文件名
% 如果图片存放在上级目录,则使用 ../../image/ 作为备用路径

% 例如:
% \begin{figure}[H]
% \centering
% \includegraphics[scale=0.3]{图.png}
% \caption{}
% \label{figure:图}
% \end{figure}
% 注意:上述\label{}一定要放在\caption{}之后,否则引用图片序号会只会显示??.

\begin{document}

\section{抽象曲面上的几何学}

\subsection{在$(M,g)$中的光滑曲线的长度}

\begin{definition}
设$\gamma:[a,b]\to M$是黎曼流形$(M,g)$上的光滑曲线，对局部容许坐标卡$(U_i,\varphi_i)$，若曲线段$\gamma|_{[t_0,t_1]}\subset U_i$，令
\begin{align*}
u_i^\alpha(t)=(\varphi_i^{-1}(\gamma(t)))^\alpha,\quad t_0\leqslant t\leqslant t_1,
\end{align*}
则$\gamma'(t)$在该坐标卡下的分量为$\left(\frac{\mathrm{d}u_i^1(t)}{\mathrm{d}t},\frac{\mathrm{d}u_i^2(t)}{\mathrm{d}t}\right)$，满足切向量分量的坐标变换规律。同时$\gamma'(t)$可视为微分算子：对$f\in C^\infty_{\gamma(t)}$，
\begin{align*}
(\gamma'(t))(f)=\frac{\mathrm{d}(f\circ\gamma(t))}{\mathrm{d}t},
\end{align*}
称$\gamma'(t)$为曲线$\gamma$的\textbf{切向量}。
\end{definition}

\begin{definition}
光滑曲线$\gamma:[a,b]\to M$的\textbf{长度}定义为
\begin{align}
L(\gamma)=\int_{a}^{b}|\gamma'(t)|\mathrm{d}t=\int_{a}^{b}\sqrt{g(\gamma'(t),\gamma'(t))}\mathrm{d}t,\label{eq:curve_length}
\end{align}
其中$|\gamma'(t)|=\sqrt{g(\gamma'(t),\gamma'(t))}$为切向量的模长。
\end{definition}

\begin{theorem}
若$t=t(s)$为保持定向的参数变换（即$t'(s)>0$），$\tilde{\gamma}(s)=\gamma(t(s))$，则$L(\tilde{\gamma})=L(\gamma)$。
\end{theorem}

\begin{proof}

由微分算子性质，$\tilde{\gamma}'(s)=\frac{\mathrm{d}t(s)}{\mathrm{d}s}\cdot\gamma'(t)$，故
\[
|\tilde{\gamma}'(s)|=t'(s)|\gamma'(t)|.
\]
因此
\[
\int_{c}^{d}|\tilde{\gamma}'(s)|\mathrm{d}s=\int_{c}^{d}|\gamma'(t)|t'(s)\mathrm{d}s=\int_{a}^{b}|\gamma'(t)|\mathrm{d}t,
\]
即曲线长度与保持定向的参数变换无关。

\end{proof}

\begin{definition}
对光滑曲线$\gamma:[a,b]\to M$，定义\textbf{弧长参数}
\begin{align*}
s=s(t)=\int_{a}^{t}|\gamma'(\tau)|\mathrm{d}\tau,
\end{align*}
满足$0\leqslant s\leqslant L(\gamma)$，$s'(t)=|\gamma'(t)|>0$；以$s$为参数时$|\tilde{\gamma}'(s)|=1$，且$L(\gamma|_{[a,t]})=s$。
\end{definition}

\subsection{$(M,g)$的一个紧致闭区域的面积}

\begin{definition}
设$X$为拓扑空间，子集$A\subset X$，若$A$的任意开覆盖均存在有限子覆盖，则称$A$为\textbf{紧致子集}；$X$的连通开子集的闭包称为$X$的\textbf{闭区域}。
\end{definition}

\begin{definition}
设$M$为二维光滑流形，若存在容许坐标卡集$\{(U_i,\varphi_i):i\in\tilde{I}\}$满足$\bigcup_{i\in\tilde{I}}U_i=M$，且$U_i\cap U_j\neq\varnothing $时坐标变换的Jacobi行列式$\frac{\partial(u_i^1,u_i^2)}{\partial(u_j^1,u_j^2)}>0$，则称$M$为\textbf{可定向流形}，该坐标卡集给出$M$的一个定向。
\end{definition}

\begin{definition}
设$D$为可定向二维黎曼流形$(M,g)$的紧致闭区域。若$D\subset U_i$（$(U_i,\varphi_i)$为定向的容许坐标卡），定义局部面积
\begin{align}
A(D)=\int_{\varphi_i^{-1}(D)}\sqrt{g_{11}^{(i)}g_{22}^{(i)}-(g_{12}^{(i)})^2}\mathrm{d}u_i^1\mathrm{d}u_i^2.\label{eq:local_area}
\end{align}
将$D$分解为有限个落在定向坐标卡内的子区域$D=\bigcup_{a=1}^N D_a$，定义整体面积
\begin{align*}
A(D)=\sum_{a=1}^N A(D_a),
\end{align*}
称为闭区域$D$的\textbf{面积}。
\end{definition}

\subsection{$(M,g)$上的切向量场的协变微分}

\begin{theorem}
局部面积\eqref{eq:local_area}的值与包含$D$的定向容许坐标卡的选取无关。
\end{theorem}

\begin{proof}

设$(U_i,\varphi_i),(U_j,\varphi_j)$为两个定向坐标卡，$D\subset U_i\cap U_j$，则度量矩阵满足
\[
\begin{pmatrix}g_{11}^{(i)}&g_{12}^{(i)}\\g_{21}^{(i)}&g_{22}^{(i)}\end{pmatrix}
=\begin{pmatrix}\frac{\partial u_j^1}{\partial u_i^1}&\frac{\partial u_j^2}{\partial u_i^1}\\\frac{\partial u_j^1}{\partial u_i^2}&\frac{\partial u_j^2}{\partial u_i^2}\end{pmatrix}
\begin{pmatrix}g_{11}^{(j)}&g_{12}^{(j)}\\g_{21}^{(j)}&g_{22}^{(j)}\end{pmatrix}
\begin{pmatrix}\frac{\partial u_j^1}{\partial u_i^1}&\frac{\partial u_j^1}{\partial u_i^2}\\\frac{\partial u_j^2}{\partial u_i^1}&\frac{\partial u_j^2}{\partial u_i^2}\end{pmatrix}.
\]
因此
\[
g_{11}^{(i)}g_{22}^{(i)}-(g_{12}^{(i)})^2=\big(g_{11}^{(j)}g_{22}^{(j)}-(g_{12}^{(j)})^2\big)\cdot\left(\frac{\partial(u_j^1,u_j^2)}{\partial(u_i^1,u_i^2)}\right)^2,
\]
\[
\sqrt{g_{11}^{(i)}g_{22}^{(i)}-(g_{12}^{(i)})^2}=\sqrt{g_{11}^{(j)}g_{22}^{(j)}-(g_{12}^{(j)})^2}\cdot\frac{\partial(u_j^1,u_j^2)}{\partial(u_i^1,u_i^2)}.
\]
由二重积分变量替换公式
\[
\int_{\varphi_j^{-1}(D)}\sqrt{g_{11}^{(j)}g_{22}^{(j)}-(g_{12}^{(j)})^2}\mathrm{d}u_j^1\mathrm{d}u_j^2
=\int_{\varphi_i^{-1}(D)}\sqrt{g_{11}^{(i)}g_{22}^{(i)}-(g_{12}^{(i)})^2}\mathrm{d}u_i^1\mathrm{d}u_i^2,
\]
故局部面积与坐标卡选取无关。

\end{proof}

\begin{definition}
设$M$为二维光滑流形，对光滑切向量场$\boldsymbol{X},\boldsymbol{Y},\boldsymbol{Z}$，$f\in C^\infty(M)$，若映射$\mathrm{D}_{\boldsymbol{X}}\boldsymbol{Y}$满足
\begin{enumerate}[(1)]
\item $\mathrm{D}_{f\boldsymbol{X}}\boldsymbol{Y}=f\mathrm{D}_{\boldsymbol{X}}\boldsymbol{Y},\quad \mathrm{D}_{\boldsymbol{X}+\boldsymbol{Z}}\boldsymbol{Y}=\mathrm{D}_{\boldsymbol{X}}\boldsymbol{Y}+\mathrm{D}_{\boldsymbol{Z}}\boldsymbol{Y}$；
\item $\mathrm{D}_{\boldsymbol{X}}(f\boldsymbol{Y})=\boldsymbol{X}(f)\boldsymbol{Y}+f\mathrm{D}_{\boldsymbol{X}}\boldsymbol{Y},\quad \mathrm{D}_{\boldsymbol{X}}(\boldsymbol{Y}+\boldsymbol{Z})=\mathrm{D}_{\boldsymbol{X}}\boldsymbol{Y}+\mathrm{D}_{\boldsymbol{X}}\boldsymbol{Z}$；
\item $\boldsymbol{X}(g(\boldsymbol{Y},\boldsymbol{Z}))=g(\mathrm{D}_{\boldsymbol{X}}\boldsymbol{Y},\boldsymbol{Z})+g(\boldsymbol{Y},\mathrm{D}_{\boldsymbol{X}}\boldsymbol{Z})$,
\end{enumerate}
则称$\mathrm{D}_{\boldsymbol{X}}\boldsymbol{Y}$为$\boldsymbol{Y}$沿$\boldsymbol{X}$的\textbf{协变导数}。
\end{definition}

\begin{definition}
在容许坐标卡$(U,\varphi)$下，设$\mathrm{D}_{\frac{\partial}{\partial u^\alpha}}\frac{\partial}{\partial u^\beta}=\Gamma_{\beta\alpha}^\gamma\frac{\partial}{\partial u^\gamma}$，其中$\Gamma_{\alpha\beta}^\gamma$为\textbf{Christoffel记号}，表达式为
\begin{align}
\Gamma_{\alpha\beta}^\gamma=\frac{1}{2}g^{\delta\gamma}\left(\frac{\partial g_{\alpha\delta}}{\partial u^\beta}+\frac{\partial g_{\delta\beta}}{\partial u^\alpha}-\frac{\partial g_{\alpha\beta}}{\partial u^\delta}\right),\label{eq:::IO235322JIOWAIJ2.}
\end{align}
式中$(g^{\alpha\beta})$为度量矩阵$(g_{\alpha\beta})$的逆矩阵，满足$g^{\alpha\gamma}g_{\gamma\beta}=\delta_\beta^\alpha$。
\end{definition}

\begin{theorem}
在局部坐标$(U;u^\alpha)$下，若$\boldsymbol{X}|_U=X^\alpha\frac{\partial}{\partial u^\alpha}$，$\boldsymbol{Y}|_U=Y^\beta\frac{\partial}{\partial u^\beta}$，则协变导数
\begin{align}
\mathrm{D}_{\boldsymbol{X}}\boldsymbol{Y}|_U=X^\alpha\left(\frac{\partial Y^\gamma}{\partial u^\alpha}+Y^\beta\Gamma_{\beta\alpha}^\gamma\right)\frac{\partial}{\partial u^\gamma},\label{eq:cov_der_local}
\end{align}
该表达式在坐标变换下不变，即$\mathrm{D}_{\boldsymbol{X}}\boldsymbol{Y}$为大范围定义在$M$上的切向量场。
\end{theorem}

\begin{proof}

由协变导数的公理化条件(1)(2)，
\[
\mathrm{D}_{\boldsymbol{X}}\boldsymbol{Y}=\mathrm{D}_{X^\alpha\frac{\partial}{\partial u^\alpha}}\left(Y^\beta\frac{\partial}{\partial u^\beta}\right)=X^\alpha\left(\frac{\partial Y^\beta}{\partial u^\alpha}\frac{\partial}{\partial u^\beta}+Y^\beta\mathrm{D}_{\frac{\partial}{\partial u^\alpha}}\frac{\partial}{\partial u^\beta}\right).
\]
代入Christoffel记号定义$\mathrm{D}_{\frac{\partial}{\partial u^\alpha}}\frac{\partial}{\partial u^\beta}=\Gamma_{\beta\alpha}^\gamma\frac{\partial}{\partial u^\gamma}$，得式\eqref{eq:cov_der_local}。
通过Christoffel记号的坐标变换公式，可验证该表达式在坐标变换下保持不变，故$\mathrm{D}_{\boldsymbol{X}}\boldsymbol{Y}$是$M$上整体定义的切向量场。

\end{proof}

\begin{definition}
对切向量场$\boldsymbol{Y}$，定义其\textbf{协变微分}
\begin{align}
\mathrm{D}\boldsymbol{Y}=\left(\mathrm{d}Y^\beta+Y^\gamma\Gamma_{\gamma\alpha}^\beta\mathrm{d}u^\alpha\right)\frac{\partial}{\partial u^\beta}.\label{eq:cov_diff_abstract}
\end{align}
\end{definition}

\subsection{切向量沿光滑曲线的平行移动}

\begin{definition}
设$\gamma:[a,b]\to M$为光滑曲线，$\boldsymbol{X}(t)$为沿$\gamma$的光滑切向量场，若
\begin{align}
\mathrm{D}_{\gamma'(t)}\boldsymbol{X}(t)=0,\label{eq:parallel_def_abstract}
\end{align}
则称$\boldsymbol{X}(t)$沿曲线$\gamma$是\textbf{平行}的。
\end{definition}

\begin{theorem}
\begin{enumerate}[(1)]
\item 切向量场$\boldsymbol{X}(t)=X^\alpha(t)\frac{\partial}{\partial u^\alpha}$沿$\gamma$平行的充要条件是
\begin{align}
\frac{\mathrm{d}X^\gamma(t)}{\mathrm{d}t}+X^\alpha(t)\frac{\mathrm{d}u^\beta(t)}{\mathrm{d}t}\Gamma_{\alpha\beta}^\gamma(\gamma(t))=0,\quad \alpha=1,2;\label{eq:parallel_ode_abstract}
\end{align}
\item 沿曲线平行的切向量场长度为常数。
\end{enumerate}
\end{theorem}

\begin{proof}

(1) 由协变导数的局部表达式，
\[
\mathrm{D}_{\gamma'(t)}\boldsymbol{X}(t)=\left(\frac{\mathrm{d}X^\gamma(t)}{\mathrm{d}t}+X^\alpha(t)\frac{\mathrm{d}u^\beta(t)}{\mathrm{d}t}\Gamma_{\alpha\beta}^\gamma(\gamma(t))\right)\frac{\partial}{\partial u^\gamma},
\]
故$\mathrm{D}_{\gamma'(t)}\boldsymbol{X}(t)=0$等价于式\eqref{eq:parallel_ode_abstract}。
(2) 由协变导数的条件(3)，
\[
\frac{\mathrm{d}}{\mathrm{d}t}\big(g(\boldsymbol{X}(t),\boldsymbol{X}(t))\big)=g(\mathrm{D}_{\gamma'(t)}\boldsymbol{X}(t),\boldsymbol{X}(t))+g(\boldsymbol{X}(t),\mathrm{D}_{\gamma'(t)}\boldsymbol{X}(t))=0,
\]
故$g(\boldsymbol{X}(t),\boldsymbol{X}(t))$为常数，即切向量长度不变。

\end{proof}

\subsection{$(M,g)$中曲线的测地曲率}

\begin{definition}
设$\gamma:[0,L]\to M$是有向二维黎曼流形$(M,g)$上以弧长$s$为参数的光滑曲线，满足$g(\gamma'(s),\gamma'(s))=1$。取$e_1(s)=\gamma'(s)$，$e_2(s)$为$e_1(s)$按曲面正定向旋转$90^\circ$得到的单位切向量，若
\begin{align}
\mathrm{D}_{\gamma'(s)}\gamma'(s)=\kappa_g(s)e_2(s),\label{eq:7.16}
\end{align}
则称
\begin{align*}
\kappa_g(s)=g\left(\mathrm{D}_{\gamma'(s)}\gamma'(s),e_2(s)\right)
\end{align*}
为曲线$\gamma(s)$的\textbf{测地曲率}。
\end{definition}

\begin{proposition}
设曲线$\gamma(s)$在容许坐标卡$(U,\varphi)$内的参数方程为$u^\alpha=u^\alpha(s)\ (\alpha=1,2)$，度量矩阵为$(g_{\alpha\beta})$，克里斯托费尔符号为$\Gamma^\alpha_{\beta\gamma}$，则测地曲率的坐标表达式为
\begin{align}
\kappa_g(s)=\sqrt{g_{11}g_{22}-(g_{12})^2}\cdot
\begin{vmatrix}
\frac{\mathrm{d}u^1}{\mathrm{d}s} & \frac{\mathrm{d}^2u^1(s)}{\mathrm{d}s^2}+\frac{\mathrm{d}u^\beta(s)}{\mathrm{d}s}\frac{\mathrm{d}u^\gamma(s)}{\mathrm{d}s}\Gamma^1_{\beta\gamma} \\
\frac{\mathrm{d}u^2}{\mathrm{d}s} & \frac{\mathrm{d}^2u^2(s)}{\mathrm{d}s^2}+\frac{\mathrm{d}u^\beta(s)}{\mathrm{d}s}\frac{\mathrm{d}u^\gamma(s)}{\mathrm{d}s}\Gamma^2_{\beta\gamma}
\end{vmatrix}.\label{eq:7.19}
\end{align}
\end{proposition}

\begin{proof}

设$\gamma'(s)=\frac{\mathrm{d}u^\alpha(s)}{\mathrm{d}s}\frac{\partial}{\partial u^\alpha}$，$e_2(s)=v^\alpha(s)\frac{\partial}{\partial u^\alpha}$，由$e_2$的单位正交性得
\begin{align}
g_{\alpha\beta}v^\alpha v^\beta=1,\quad g_{\alpha\beta}v^\alpha\frac{\mathrm{d}u^\beta(s)}{\mathrm{d}s}=0.\label{eq:7.17}
\end{align}
由\eqref{eq:7.17}第二式得$v^\alpha g_{\alpha1}\frac{\mathrm{d}u^1}{\mathrm{d}s}+v^\alpha g_{\alpha2}\frac{\mathrm{d}u^2}{\mathrm{d}s}=0$，故
\begin{align}
\left(v^\alpha g_{\alpha1},v^\alpha g_{\alpha2}\right)=\lambda\left(-\frac{\mathrm{d}u^2}{\mathrm{d}s},\frac{\mathrm{d}u^1}{\mathrm{d}s}\right),\quad \lambda>0.\label{eq:7.18}
\end{align}
由黎曼联络的坐标公式
\begin{align*}
\mathrm{D}_{\gamma'(s)}\gamma'(s)=\left(\frac{\mathrm{d}^2u^\alpha(s)}{\mathrm{d}s^2}+\frac{\mathrm{d}u^\beta(s)}{\mathrm{d}s}\frac{\mathrm{d}u^\gamma(s)}{\mathrm{d}s}\Gamma^\alpha_{\beta\gamma}(\gamma(s))\right)\left.\frac{\partial}{\partial u^\alpha}\right|_{\gamma(s)},
\end{align*}
代入\eqref{eq:7.16}，结合\eqref{eq:7.18}化简得
\begin{align*}
\kappa_g(s)=\lambda\cdot
\begin{vmatrix}
\frac{\mathrm{d}u^1}{\mathrm{d}s} & \frac{\mathrm{d}^2u^1(s)}{\mathrm{d}s^2}+\frac{\mathrm{d}u^\beta(s)}{\mathrm{d}s}\frac{\mathrm{d}u^\gamma(s)}{\mathrm{d}s}\Gamma^1_{\beta\gamma} \\
\frac{\mathrm{d}u^2}{\mathrm{d}s} & \frac{\mathrm{d}^2u^2(s)}{\mathrm{d}s^2}+\frac{\mathrm{d}u^\beta(s)}{\mathrm{d}s}\frac{\mathrm{d}u^\gamma(s)}{\mathrm{d}s}\Gamma^2_{\beta\gamma}
\end{vmatrix}.
\end{align*}
将\eqref{eq:7.18}改写为矩阵形式
\begin{align*}
\left(v^1,v^2\right)\begin{pmatrix}g_{11}&g_{12}\\g_{21}&g_{22}\end{pmatrix}=\lambda\left(-\frac{\mathrm{d}u^2}{\mathrm{d}s},\frac{\mathrm{d}u^1}{\mathrm{d}s}\right),
\end{align*}
解得
\begin{align*}
\left(v^1,v^2\right)=\frac{\lambda}{g_{11}g_{22}-(g_{12})^2}\left(-g_{2\alpha}\frac{\mathrm{d}u^\alpha}{\mathrm{d}s},g_{1\alpha}\frac{\mathrm{d}u^\alpha}{\mathrm{d}s}\right).
\end{align*}
代入$g_{\alpha\beta}v^\alpha v^\beta=1$，可求得$\lambda=\sqrt{g_{11}g_{22}-(g_{12})^2}$，代入即得\eqref{eq:7.19}.

\end{proof}

\subsection{$(M,g)$中的测地线}

\begin{definition}
二维黎曼流形$(M,g)$上的\textbf{测地线}为测地曲率$\kappa_g=0$的光滑曲线，等价于其单位切向量$\gamma'(s)$沿曲线$\gamma(s)$平行移动。
\end{definition}

\begin{proposition}
测地线$\gamma(s)$的坐标参数方程$u^\alpha=u^\alpha(s)$满足微分方程组
\begin{align}
\frac{\mathrm{d}^2u^\alpha(s)}{\mathrm{d}s^2}+\frac{\mathrm{d}u^\beta(s)}{\mathrm{d}s}\frac{\mathrm{d}u^\gamma(s)}{\mathrm{d}s}\Gamma^\alpha_{\beta\gamma}(\gamma(s))=0,\quad \alpha=1,2.\label{eq:geo_eq}
\end{align}
\end{proposition}
\begin{proof}

由测地线定义$\kappa_g=0$，即$\mathrm{D}_{\gamma'(s)}\gamma'(s)=0$，代入黎曼联络的坐标表达式直接可得。

\end{proof}



\end{document}